\title{DeepSeekMath: Pushing the Limits of Mathematical Reasoning in Open Language Models}

\begin{abstract}
Mathematical reasoning remains a formidable challenge for language models, given its inherently complex and highly structured characteristics. In this work, we present DeepSeekMath~7B, an extension of DeepSeek-Coder-Base-v1.5 7B that is further pre-trained on a substantial 120B math-focused tokens extracted from Common Crawl, supplemented with both natural language and programming code datasets. DeepSeekMath~7B attains a notable score of 51.7\% on the MATH competition-level benchmark, achieving this without the aid of external toolkits or ensembling strategies, and reaching performance comparable to Gemini-Ultra and GPT-4. When employing self-consistency over 64 generations, DeepSeekMath~7B improves to 60.9\% accuracy on the MATH benchmark. The superior mathematical reasoning performance of DeepSeekMath~can be credited to two principal innovations: Firstly, we exploit the extensive resources of web-based data through a carefully designed data selection methodology. Secondly, we propose Group Relative Policy Optimization (GRPO), a new approach derived from Proximal Policy Optimization (PPO), which both elevates mathematical problem-solving abilities and enhances the memory efficiency of PPO during training.
\end{abstract}

\section{Introduction}

Large language models (LLM) have fundamentally transformed the field of mathematical reasoning within artificial intelligence, driving remarkable progress in quantitative reasoning benchmarks \citep{MATH} as well as geometry reasoning assessments \citep{trinh2024solving}. Furthermore, such models have demonstrated value in supporting human users as they tackle intricate mathematical challenges \citep{tao}. Despite these advances, state-of-the-art systems such as GPT-4 \citep{gpt4} and Gemini-Ultra \citep{gemini} remain closed, and the available open-source alternatives exhibit a marked gap in performance relative to these proprietary models.

This paper presents DeepSeekMath, a specialized language model focused on mathematics that markedly surpasses existing open-source models in mathematical reasoning and approaches the proficiency of GPT-4 on established academic assessments. Central to this advancement is the construction of the DeepSeekMath Corpus, an extensive, high-quality pre-training dataset encompassing 120B math-related tokens. The corpus is curated from Common Crawl (CC) through a classifier built on fastText \citep{joulin2016fasttext}. Initially, the classifier is trained with examples from OpenWebMath \citep{openwebmath} as positives and a varied array of unrelated web content as negatives. The trained classifier is then applied to the CC to discover further math-relevant content, which undergoes human validation and refinement. After incorporating this human-annotated data, the classifier is retrained to enhance its accuracy. Evaluation demonstrates the effectiveness of this large-scale corpus, with DeepSeekMath-Base 7B achieving 64.2\% on GSM8K \citep{gsm8k} and 36.2\% on the MATH competition benchmark \citep{MATH}, surpassing the performance of Minerva 540B \citep{minerva}. Furthermore, since the DeepSeekMath Corpus includes multiple languages, we observe improved outcomes on Chinese mathematics benchmarks \citep{wei2023cmath,agieval}. We view our mathematical data processing pipeline as a useful contribution for further research and anticipate that future work can achieve even greater advancements.

DeepSeekMath-Base is built upon the DeepSeek-Coder-Base-v1.5 7B model \citep{deepseek-coder}, as our findings suggest that initializing from a code-focused model yields better results than leveraging a general large language model. Additionally, we find that mathematical pre-training enhances not only the model’s proficiency in mathematics but also its performance on broader reasoning assessments, such as MMLU \citep{mmlu} and BBH \citep{bbh}, demonstrating an improvement in both mathematical and general reasoning capacities.

Following pre-training, we perform instruction tuning on DeepSeekMath-Base using datasets that incorporate chain-of-thought reasoning \citep{cot}, program-of-thought frameworks \citep{pot,pal}, and tool-integrated problem-solving approaches \citep{tora}. The tuned model, DeepSeekMath-Instruct 7B, outperforms all other models in the 7B category and is competitive with instruction-tuned models containing up to 70B parameters in the open-source domain.

In addition, we propose Group Relative Policy Optimization (GRPO), a modified reinforcement learning algorithm derived from Proximal Policy Optimization (PPO) \citep{schulman2017proximal}. GRPO departs from traditional approaches by eliminating the need for a critic network, opting instead to estimate the learning baseline from group-level scores. This change considerably decreases the computational demand for training. Utilizing only a select portion of English instruction tuning data, GRPO delivers marked advancements over the robust DeepSeekMath-Instruct baseline, with notable gains observed in both in-domain (GSM8K: 82.9\% $\rightarrow$ 88.2\%, MATH: 46.8\% $\rightarrow$ 51.7\%) and out-of-domain mathematical tasks (such as CMATH: 84.6\% $\rightarrow$ 88.8\%) during the reinforcement phase. Moreover, we propose a unified framework for understanding various methodologies, including Rejection Sampling Fine-Tuning (RFT) \citep{yuan2023scaling}, Direct Preference Optimization (DPO) \citep{dpo}, PPO, and GRPO. Within this framework, these approaches are interpreted as forms of direct or simplified reinforcement learning. To further elucidate the core elements of this paradigm, we conduct thorough empirical studies comparing strategies such as online versus offline learning, supervision based on outcomes versus processes, and single-step versus iterative reinforcement learning. Finally, we offer insights into the mechanisms by which reinforcement learning augments instruction-tuned model performance, and outline promising avenues for advancing RL methods within this comprehensive framework.

\subsection{Contributions}
This work advances large-scale math pre-training and investigates reinforcement learning techniques through a thorough analytical approach.

\textbf{Math Pre-Training at Scale}

\begin{itemize}[topsep=0pt]
    \item We demonstrate that the Common Crawl dataset, which is freely available, possesses substantial value for mathematical applications. Employing a carefully crafted data filtration pipeline, we assembled the DeepSeekMath Corpus—a meticulously curated collection containing 120B tokens sourced from web pages specifically selected for mathematical relevance. This resource is nearly 7 times greater than the math web data utilized by Minerva \citep{minerva} and exceeds the recent OpenWebMath corpus \citep{openwebmath} by a factor of 9.

    \item Our DeepSeekMath-Base 7B model, trained from scratch, achieves results on par with Minerva 540B \citep{minerva}. This demonstrates that a model's parameter count is not the sole determinant of proficiency in mathematical reasoning. Robust outcomes can be obtained from more compact models if pre-trained with superior-quality data.

    \item We report insights from our math pre-training experiments. Notably, performing code pre-training prior to math training boosts a model’s proficiency in tackling mathematical tasks, both with and without auxiliary tool support. This provides evidence toward addressing the enduring debate: \textit{does code training enhance reasoning capabilities?} Our results support that, at least regarding mathematical reasoning, code training is beneficial.

    \item While utilizing arXiv papers as training material is a prevalent practice, especially within the mathematical domain, our experiments reveal that such data offers no substantial benefits across any of the mathematical benchmarks assessed in this study.
\end{itemize}

\textbf{Exploration and Analysis of Reinforcement Learning}

\begin{itemize}[topsep=0pt]
    \item We present Group Relative Policy Optimization (GRPO), a novel reinforcement learning approach that operates without a critic model, instead leveraging group scores to estimate the baseline. This strategy substantially reduces the computational demands associated with training in comparison to Proximal Policy Optimization (PPO).
    \item We show that GRPO notably improves the capabilities of our instruction-tuned model, DeepSeekMath-Instruct, by using only instruction-tuning data. In addition, we observe gains in out-of-domain task performance throughout the reinforcement learning phase.
    \item We propose a cohesive framework to analyze various techniques, such as RFT, DPO, PPO, and GRPO. Through extensive empirical study—including comparisons of online versus offline training, outcome versus process supervision, and single-step versus iterative reinforcement learning—we delve into the fundamental aspects underlying this framework.
    \item Within this unified perspective, we examine the mechanisms contributing to the effectiveness of reinforcement learning, and outline several promising pathways for advancing reinforcement learning with large language models (LLMs).
\end{itemize}

\subsection{Summary of Evaluations and Metrics}

\begin{itemize}[topsep=0pt]
    \item \textbf{English and Chinese Mathematical Reasoning}: We thoroughly evaluate our models on a range of English and Chinese benchmarks, spanning mathematical questions from elementary school up to college difficulty. The English evaluation suite consists of GSM8K \citep{gsm8k}, MATH \citep{MATH}, SAT \citep{llemma}, OCW Courses \citep{minerva}, and MMLU-STEM \citep{mmlu}, while the Chinese set includes MGSM-zh \citep{mgsm}, CMATH \citep{wei2023cmath}, Gaokao-MathCloze \citep{agieval}, and Gaokao-MathQA \citep{agieval}. Our assessment focuses on each model's proficiency in generating comprehensive, stand-alone written solutions, both without reliance on external tools and when permitted to use Python.

    In the English benchmarks, DeepSeekMath-Base achieves results comparable to the proprietary Minerva 540B \citep{minerva}, and outperforms every available open-source base model—such as Mistral 7B \citep{mistral} and Llemma-34B \citep{llemma}—regardless of their math pre-training regimen, often by a considerable difference. Importantly, DeepSeekMath-Base excels on the Chinese benchmarks, which may be attributed to our inclusion of high-quality non-English math data, in contrast to prior studies \citep{minerva,llemma} that restricted pre-training to English datasets. Following mathematical instruction tuning and reinforcement learning, DeepSeekMath-Instruct and DeepSeekMath-RL both demonstrate robust outcomes, achieving—for the first time in the open-source arena—accuracy surpassing 50\% on the advanced-level MATH benchmark.

\item \textbf{Formal Mathematics}: We assess DeepSeekMath-Base’s performance on the informal-to-formal theorem proving task introduced by \citep{dsp_proof} using the miniF2F benchmark \citep{minif2f}, with Isabelle \citep{isabelle} serving as the proof assistant. DeepSeekMath-Base achieves robust results in few-shot autoformalization.

\item \textbf{Natural Language Understanding, Reasoning, and Code}: To obtain a broad evaluation of the models’ capabilities in general comprehension, reasoning, and programming, we test DeepSeekMath-Base across several benchmarks. These include the Massive Multitask Language Understanding (MMLU) benchmark \citep{mmlu}, which spans 57 multiple-choice tasks from a variety of disciplines; BIG-Bench Hard (BBH) \citep{bbh}, which comprises 23 especially challenging, mostly multi-step reasoning tasks; as well as HumanEval \citep{codex} and MBPP \citep{mbpp}, both standard evaluations for code generation models. The results suggest that pre-training on mathematical content strengthens both understanding and reasoning in language tasks.

\section{Math Pre-Training}

\subsection{Data Collection and Decontamination} In this subsection, we describe how the DeepSeekMath Corpus is assembled from Common Crawl. Figure \ref{fig:data_collect} presents a stepwise, iterative workflow, illustrating the methodology for systematically extracting a large-scale math-specific corpus from Common Crawl. The process begins with a high-quality, albeit relatively small, seed dataset focused on mathematics. It is important to highlight that this procedure is adaptable to various domains beyond mathematics, such as computer programming.

We start by selecting OpenWebMath \citep{openwebmath}—an assemblage of premium mathematical web content—as our initial seed dataset. This data is used to train a fastText model \citep{joulin2016fasttext}, enabling the retrieval of additional web pages that exhibit characteristics similar to OpenWebMath. To this end, 500,000 items are randomly sampled from the seed set as positive examples, while another 500,000 Common Crawl web pages serve as negative samples. The training utilizes an open-source toolkit\footnote{\url{https://fasttext.cc}} configured with a vector size of 256, a learning rate of 0.1, a word n-gram length maximum of 3, a minimum word occurrence threshold of 3, and 3 training epochs. To efficiently narrow down Common Crawl’s dataset, URL-based deduplication along with near-deduplication methods are applied, reducing the dataset to 40B HTML web pages. Mathematical pages are then retrieved from the deduplicated collection via the trained fastText model. To further eliminate poor-quality content, these pages are ranked by the fastText model’s scoring, and only those with the highest scores are retained. The data retention volume is determined by conducting pre-training experiments with top 40B, 80B, 120B, and 160B tokens. For the initial round, we select the top 40B tokens.

Following the initial round of data acquisition, a significant number of mathematical web pages remain uncovered. This is largely due to the limited variety within the positive examples used to train the fastText model. To address this limitation, we seek out additional mathematical web domains to augment the initial dataset, thereby facilitating further optimization of the fastText model. Our methodology begins by segmenting the entire Common Crawl dataset into non-overlapping domains, where each domain consists of web pages sharing an identical base URL. We then compute the proportion of pages collected during the first iteration for every domain. If more than 10\% of a domain’s pages are retrieved, the domain is deemed math-focused (e.g., \url{mathoverflow.net}).

For those domains identified as mathematical, we perform a manual annotation process on the URLs that correspond to math-related sections (e.g., \url{mathoverflow.net/questions}). Pages linked to these annotated URLs but not yet collected are incorporated into the seed dataset. This strategy expands our set of positive examples and permits us to refine the fastText model so that it retrieves a larger quantity of mathematical content in the subsequent round. Upon four rounds of data collection, our final dataset comprises 35.5M mathematical web pages, encompassing 120B tokens. Notably, by the fourth round, approximately 98\% of the available data had already been captured by the end of the third round, prompting the termination of further data gathering.

To mitigate the risk of benchmark leakage, we employ the filtering approach outlined in \cite{deepseek-coder}. Specifically, we exclude any web pages containing questions or answers from notable English mathematical benchmarks such as GSM8K \citep{gsm8k} and MATH \citep{MATH}, as well as Chinese resources like CMATH \citep{wei2023cmath} and AGIEval \citep{agieval}. Our filtering procedure dictates that any text snippet in the corpus containing a 10-gram sequence that matches exactly with a substring in any evaluation benchmark is discarded. For benchmark texts comprising between 3 and 9 grams, we apply an exact match to ensure the exclusion of compromised web pages from our mathematical training corpus.

\subsection{Validating the Quality of the DeepSeekMath~Corpus}

To assess the utility of the DeepSeekMath~Corpus, we conduct pre-training trials that contrast its quality against recently released corpora aimed at mathematical training:

\begin{itemize}[topsep=0pt]
    \item \textbf{MathPile} \citep{mathpile}: A composite corpus totaling 8.9B tokens sourced from multiple repositories, including textbooks, Wikipedia, ProofWiki, CommonCrawl, StackExchange, and arXiv; with arXiv contributing over 85\% of the data;
    \item \textbf{OpenWebMath} \citep{openwebmath}: Extracted from CommonCrawl, filtered for mathematical material, with a total of 13.6B tokens;
    \item \textbf{Proof-Pile-2} \citep{llemma}: A mathematical collection comprising OpenWebMath, AlgebraicStack (which contains 10.3B tokens of mathematical code), and arXiv papers (28.0B tokens). For experiments on Proof-Pile-2, we adhere to the arXiv:Web:Code ratio of 2:4:1 as specified by \cite{llemma}.
\end{itemize}

\subsubsection{Training Setting}
\label{sec:quality-policy}
We perform math-specific fine-tuning on a general language model containing 1.3B parameters, architecturally consistent with the DeepSeek LLMs \citep{deepseek-llm} and referred to as DeepSeek-LLM 1.3B. Separate training runs are conducted for each mathematical dataset, each for 150B tokens. All model training is carried out using the lightweight and efficient HAI-LLM \citep{haillm} framework. Adhering to the training methodology adopted by DeepSeek LLMs, we employ the AdamW optimizer \citep{adamW} with hyperparameters set to $\beta_1=0.9$, $\beta_2=0.95$, and $\mathrm{weight\_decay}=0.1$. The learning rate follows a multi-stage schedule: it increases to its peak after 2,000 warmup iterations, declines to 31.6\% of its maximum value after 80\% of the total training steps, and then further drops to 10.0\% of the peak after 90\% completion. The peak learning rate is established at 5.3e-4. Training uses a batch size of 4 million tokens and a context window of 4,000 tokens.

\subsubsection{Evaluation Results}

\textbf{DeepSeekMath~Corpus stands out for its superior quality, multilingual coverage, and unmatched scale.}

\begin{itemize}[topsep=0pt]
    \item \textbf{High-quality}:
    We assess performance on eight mathematical evaluation suites using a few-shot chain-of-thought prompting protocol \cite{cot}. Table \ref{tab:corpora_comparison} clearly illustrates that models fine-tuned on DeepSeekMath~Corpus consistently achieve higher scores. Furthermore, Figure \ref{fig:corpora_comparisons} demonstrates that, at the 50B token mark (equivalent to one epoch on Proof-Pile-2), the DeepSeekMath-trained model outperforms the Proof-Pile-2 counterpart, suggesting that DeepSeekMath~Corpus provides overall better-quality data.
    \item \textbf{Multilingual}:
    DeepSeekMath~Corpus incorporates texts from a range of languages, with English and Chinese as the principal contributors. As indicated in Table \ref{tab:corpora_comparison}, models trained on DeepSeekMath~Corpus exhibit enhanced mathematical reasoning abilities in both English and Chinese. This contrasts with previous mathematical corpora, which are predominantly English-focused and provide little to no improvement—and occasionally negative impact—on mathematical reasoning tasks in Chinese.
    \item \textbf{Large-scale}:
    DeepSeekMath~Corpus significantly surpasses the size of previously released mathematical datasets. Figure \ref{fig:corpora_comparisons} reveals that DeepSeek-LLM 1.3B, when trained on DeepSeekMath~Corpus, displays a much more pronounced learning curve with sustained performance gains. By comparison, baseline corpora are much smaller, necessitating multiple passes during training, which causes the resulting model improvements to plateau quickly.
\end{itemize}

\subsection{Training and Evaluating DeepSeekMath-Base 7B}

Here, we present DeepSeekMath-Base 7B, a foundational model notable for its advanced reasoning skills, particularly within the domain of mathematics. The model's initial parameters are inherited from DeepSeek-Coder-Base-v1.5 7B \citep{deepseek-coder}, upon which we conduct further training using a dataset of 500B tokens. The composition of this training data is as follows: 56\% is drawn from the DeepSeekMath Corpus, 4\% from AlgebraicStack, 10\% from arXiv, 20\% from Github code repositories, and the final 10\% comprises natural language material sourced from Common Crawl, available in both English and Chinese. For training configuration, we primarily adhere to the protocol outlined in Section \ref{sec:quality-policy}, with the exception of setting the peak learning rate to 4.2e-4 and employing a batch size of 10M tokens.

We undertake an extensive evaluation of DeepSeekMath-Base 7B's proficiency in mathematics, examining the model's capabilities in independently generating coherent mathematical solutions, utilizing tools for problem solving, and executing formal theorem proving tasks. In addition to mathematical assessment, we report on the model’s broader capabilities, such as natural language comprehension, logical reasoning, and programming aptitude.

\paragraph{Mathematical Problem Solving with Step-by-Step Reasoning}

To assess DeepSeekMath-Base's ability to tackle mathematical problems, we utilize few-shot chain-of-thought prompting \citep{cot} across eight different evaluation suites in both English and Chinese. These benchmarks include those centered on quantitative reasoning (such as GSM8K \citep{gsm8k}, MATH \citep{MATH}, and CMATH \citep{wei2023cmath}), as well as multiple-choice assessments (including MMLU-STEM \citep{mmlu} and Gaokao-MathQA \citep{agieval}). Together, these benchmarks span a wide array of mathematical topics and difficulty levels, ranging from elementary school concepts to college-level problems.

As detailed in Table \ref{tab:base_math_cot}, DeepSeekMath-Base 7B achieves leading results across all eight evaluated benchmarks relative to open-source base models, such as the widely-adopted generalist Mistral 7B \citep{mistral} and the more recent Llemma 34B \citep{llemma}, which received mathematical pretraining using Proof-Pile-2 \citep{llemma}. Remarkably, on the MATH dataset, which tests competition-level mathematical ability, DeepSeekMath-Base exceeds all previously available open-source base models by more than 10\% in absolute terms. It also outperforms the closed-source Minerva 540B \citep{minerva}—a model 77 times its size, based on PaLM \citep{palm} and further pretrained on mathematics-focused corpora.

\paragraph{Mathematical Problem Solving with Tool Use} For evaluating program-aided mathematical reasoning, we applied few-shot program-of-thought prompting \citep{pot,pal} on GSM8K and MATH. In this setup, models are asked to resolve each problem by authoring a Python script that may use libraries like \textit{math} and \textit{sympy} to handle complex calculations; the result of script execution serves as the predicted solution. According to Table \ref{tab:base_math_pot_proof}, DeepSeekMath-Base 7B establishes a new state-of-the-art among open-source models, surpassing Llemma 34B on these tasks.

\paragraph{Formal Mathematics}

The automation of formal proofs is increasingly recognized for bolstering both rigor and productivity in mathematics. We test DeepSeekMath-Base 7B on the informal-to-formal proof generation task from \citep{dsp_proof}, which requires transforming an informal statement, its formal version, and an informal proof into a formalized proof. Evaluations take place on miniF2F \citep{minif2f}—a benchmark for Olympiad-level mathematics in formal settings—where the model generates Isabelle proofs for each example via few-shot prompts. Following \cite{dsp_proof}, we employ models to construct proof outlines, after which the Sledgehammer automated prover \citep{sledgehammer} completes the formalization by resolving any outstanding gaps. As reflected in Table \ref{tab:base_math_pot_proof}, DeepSeekMath-Base 7B exhibits robust capabilities in automatic formal proof generation.

\paragraph{Natural Language Understanding, Reasoning, and Code}

To assess broader model ability, we consider performance on MMLU \citep{mmlu} for language understanding, BBH \citep{bbh} for reasoning, and HumanEval \citep{codex} alongside MBPP \citep{mbpp} for code generation. Table \ref{tab:understanding_reasoning_code} indicates that DeepSeekMath-Base 7B delivers considerable improvements on MMLU and BBH compared to its predecessor, DeepSeek-Coder-Base-v1.5 \citep{deepseek-coder}, underscoring the value of mathematics-oriented training for language understanding and reasoning tasks. Furthermore, by integrating code tokens during continual pretraining, DeepSeekMath-Base 7B maintains the coding proficiency of DeepSeek-Coder-Base-v1.5 on both code-related benchmarks. Collectively, DeepSeekMath-Base 7B demonstrates clear superiority over Mistral 7B \citep{mistral} across the three reasoning and coding evaluation datasets.

\section{Supervised Fine-Tuning}

\subsection{SFT Data Curation}

We compile an instruction-tuning dataset focused on mathematics that encompasses both English and Chinese problems across a spectrum of mathematical topics and difficulty levels. Each problem is accompanied by a solution in chain-of-thought (CoT) \citep{cot}, program-of-thought (PoT) \citep{pot,pal}, or a tool-integrated reasoning style \citep{tora}. The full training corpus comprises 776K examples.

\begin{itemize}[topsep=0pt]
    \item \textbf{English mathematical datasets}: Our curation annotates GSM8K and MATH questions with tool-integrated answer formats, and also leverages a portion of MathInstruct \citep{MathInstruct} and the Lila-OOD training split \citep{lila}, containing problems addressed via CoT or PoT approaches. This English dataset collection spans a broad array of mathematical subdisciplines such as algebra, probability, number theory, calculus, and geometry.
    \item \textbf{Chinese mathematical datasets}: For Chinese, we aggregate K-12 mathematics questions distributed over 76 distinct topics, such as linear equations, and enrich them with annotations in both CoT and tool-integrated reasoning formats.
\end{itemize}

\subsection{Training and Evaluating DeepSeekMath-Instruct 7B}

Here, we present the DeepSeekMath-Instruct 7B, which has undergone instruction tuning for mathematical problem-solving built upon the DeepSeekMath-Base foundation. During training, instances are concatenated at random, up to a maximum of 4K tokens per context window. The model is optimized for 500 iterations using a batch size of 256 and a constant learning rate of $5\times10^{-5}$.

We assess the mathematical reasoning abilities of our models under both tool-free and tool-utilizing settings, employing four quantitative reasoning benchmarks in both English and Chinese. Our results are compared against contemporary state-of-the-art systems:

\begin{itemize}[topsep=0pt]
    \item \textbf{Closed-source models} consist of: (1) the GPT series, notably GPT-4 \citep{gpt4} and GPT-4 Code Interpreter~\footnote{\url{https://openai.com/blog/chatgpt-plugins\#\#code-interpreter}}, (2) Gemini Ultra and Gemini Pro \citep{gemini}, (3) Inflection-2 \citep{inflection-2}, (4) Grok-1~\footnote{\url{https://x.ai/model-card}}, in addition to several newly available Chinese models, such as (5) Baichuan-3~\footnote{\url{https://www.baichuan-ai.com}}, and (6) the most recent GLM-4~\footnote{\url{https://open.bigmodel.cn/dev/api\#glm-4}} within the GLM model family \citep{glm}.
\end{itemize}

These models are designed for general applications, and the majority have been subject to various alignment processes. 
    \item \textbf{Open-source models} include: general-purpose models such as (1) DeepSeek-LLM-Chat 67B \citep{deepseek-llm}, (2) Qwen 72B \citep{qwen}, (3) SeaLLM-v2 7B \citep{seallm}, and (4) ChatGLM3 6B \citep{chatglm3}; alongside models tailored for mathematical capabilities, including (5) InternLM2-Math 20B~\footnote{\url{https://github.com/InternLM/InternLM-Math}}, which extends InternLM2 through additional math-focused pre-training and subsequent instruction tuning; (6) Math-Shepherd-Mistral 7B, utilizing PPO training \citep{schulman2017proximal} on Mistral 7B \citep{mistral} with a process-supervised reward structure; (7) the WizardMath models \citep{wizardmath}, which enhance the mathematical reasoning abilities of Mistral 7B and Llama-2 70B \citep{llama2} by means of evolve-instruct (an instruction tuning paradigm leveraging AI-generated instructions) and PPO, with their training datasets mainly comprised of GSM8K and MATH problems; (8) MetaMath 70B \citep{metamath}, which further fine-tunes Llama-2 70B on an enriched version of GSM8K and MATH; (9) ToRA 34B \cite{tora}, which adapts CodeLlama 34B for integrated mathematical reasoning with external tools; and (10) MAmmoTH 70B \citep{MathInstruct}, where Llama-2 70B is instruction-tuned using MathInstruct.

Referencing Table \ref{tab:sft_rl_math}, in the evaluation setup that restricts tool usage, DeepSeekMath-Instruct 7B exhibits notable proficiency in sequential reasoning. In particular, when evaluated on the competition-grade MATH benchmark, this model outperforms all other open-source alternatives and most closed-source systems (including Inflection-2 and Gemini Pro) by a margin of at least 9 percentage points. This outcome holds even when compared to much larger models (e.g., Qwen 72B) or those employing reinforcement learning specifically targeting mathematical capabilities (e.g., WizardMath-v1.1 7B). Although DeepSeekMath-Instruct demonstrates competitive results relative to Chinese proprietary models such as GLM-4 and Baichuan-3 on the MATH dataset, it continues to trail behind models like GPT-4 and Gemini Ultra.

In scenarios permitting the combination of natural language inference and programming tool integration, DeepSeekMath-Instruct 7B nearly reaches 60\% accuracy on the MATH dataset, achieving results superior to all other currently available open-source models. Furthermore, across alternative evaluation benchmarks, its performance is on par with DeepSeek-LLM-Chat 67B—the preceding open-source state-of-the-art—despite the latter being an order of magnitude larger in scale.

\section{Reinforcement Learning}

\subsection{Group Relative Policy Optimization} 
Reinforcement learning (RL) has demonstrated efficacy in further enhancing the mathematical reasoning performance of LLMs following Supervised Fine-Tuning (SFT) \citep{wang2023math,wizardmath}. In the following, we propose Group Relative Policy Optimization (GRPO), an RL method that is both efficient and effective.

\subsubsection{From PPO to GRPO}
Proximal Policy Optimization (PPO) \citep{schulman2017proximal} stands as a widely adopted actor-critic RL method during the RL fine-tuning phase for LLMs \citep{ouyang2022training}. This algorithm trains LLMs by maximizing the following clipped surrogate objective:

\begin{equation}
\footnotesize
    \mathcal{J}_{PPO}(\theta) = \mathbb{E}{[q \sim P(Q), o \sim \pi_{\theta_{old}}(O|q)]} \frac{1}{|o|} \sum_{t=1}^{|o|} \min \left[ \frac{\pi_\theta(o_{t} | q, o_{<t})}{\pi_{\theta_{old}}(o_{t} | q, o_{<t})} A_{t}, \text{clip} \left( \frac{\pi_\theta(o_{t} | q, o_{<t})}{\pi_{\theta_{old}}(o_{t} | q, o_{

In this context, $\pi_{\theta}$ and $\pi_{\theta_{old}}$ denote the current and previous policy models, respectively, while $q$ and $o$ represent questions and outputs, with the latter sampled from the question dataset and the old policy $\pi_{\theta_{old}}$. The parameter $\varepsilon$ is a clipping hyper-parameter that PPO introduces to enhance the stability of the training process. The advantage term, $A_t$, is computed using Generalized Advantage Estimation (GAE) \citep{gae}, which relies on the rewards $\{r_{\ge t}\}$ and a value function $V_{\psi}$ that is learned during training. Consequently, PPO necessitates training a value function concurrently with the policy model. Furthermore, to counteract the potential issue of reward model over-optimization, a widely adopted technique involves appending a KL penalty (per token) from a fixed reference model to the reward for each token \citep{ouyang2022training}, as follows:

\begin{equation}
    r_{t} = r_\varphi(q, o_{\le t}) - \beta \log\frac{\pi_{\theta}(o_{t}|q, o_{<t})}{\pi_{ref}(o_{t}|q, o_{<t})},
\label{eq:PPO-reward}
\end{equation}

where $r_\varphi$ stands for the reward model, $\pi_{ref}$ denotes the reference model (commonly the initial SFT model), and $\beta$ indicates the weight assigned to the KL penalty.

However, employing a value function in PPO often means allocating resources to a model with capacity similar to that of the policy model, thereby substantially increasing both memory usage and computational requirements. In reinforcement learning settings, the value function serves as a baseline in advantage estimation to help reduce variance. Within large language model scenarios, the reward model usually provides a reward score only for the final token, making it challenging to effectively train a value function to accurately predict rewards at every token position. To circumvent these challenges, we introduce Group Relative Policy Optimization (GRPO), illustrated in Figure \ref{fig:grpo}. GRPO dispenses with the explicit value function approximation required in PPO. Instead, it employs the mean reward over a set of sampled outputs—each generated in response to the same question—as a baseline. More precisely, for each question $q$, GRPO samples a batch of outputs $\{o_1, o_2, \cdots, o_G\}$ from the previous policy $\pi_{\theta_{old}}$ and refines the policy model by optimizing the following objective:

\begin{equation}
\footnotesize
\begin{split}
    \mathcal{J}_{GRPO}(\theta) &= \mathbb{E}{[q \sim P(Q), \{o_i\}_{i=1}^G \sim \pi_{\theta_{old}}(O|q)]}  \\
    & \frac{1}{G}\sum_{i=1}^G\frac{1}{|o_i|} \sum_{t=1}^{|o_i|} \left\{ \min \left[ \frac{\pi_\theta(o_{i,t} | q, o_{i,<t})}{\pi_{\theta_{old}}(o_{i,t} | q, o_{i,<t})} \hat{A}_{i,t}, \text{clip} \left( \frac{\pi_\theta(o_{i,t} | q, o_{i,<t})}{\pi_{\theta_{old}}(o_{i,t} | q, o_{i,<t})}, 1 - \varepsilon, 1 + \varepsilon \right)  \hat{A}_{i,t} \right] - \beta \mathbb{D}_{KL}\left[\pi_{\theta} || \pi_{ref}\right]\right\} ,
\end{split}
\label{eq:GRPO-obj}
\end{equation}

Here, $\varepsilon$ and $\beta$ are again hyper-parameters, while $\hat{A}_{i,t}$ denotes the advantage computed solely from the relative rewards among outputs within the same group—details of which will be elaborated in subsequent subsections. The groupwise, relative methodology adopted by GRPO for calculating advantages fits naturally with the reward models’ comparative structure, since these models are generally trained using pairwise or groupwise comparisons for responses to identical questions. It is also important to point out that, contrary to standard PPO, GRPO incorporates the KL divergence penalty as a direct regularization term in the loss function—rather than embedding it in the reward computation—thereby streamlining the computation of $\hat{A}_{i,t

\begin{algorithm}[t]
  \small
  \caption{Iterative Group Relative Policy Optimization}
  \textbf{Input} initial policy model $\pi_{\theta_{\text{init}}}$; reward models $r_\varphi$; task prompts $\mathcal{D}$; 
  hyperparameters $\varepsilon$, $\beta$, $\mu$
  \begin{algorithmic}[1]
    \State Set policy model $\pi_\theta \leftarrow \pi_{\theta_{\text{init}}}$
    \For{iteration = 1, \dots, I}
       \State Assign reference model $\pi_{ref} \leftarrow \pi_{\theta}$
      \For{step = 1, \dots, M}
      \State Draw a batch $\mathcal{D}_b$ from $\mathcal{D}$
      \State Update the old policy: $\pi_{\theta_{old}} \leftarrow \pi_{\theta}$ 
      \State For every question $q \in \mathcal{D}_b$, sample $G$ outputs $\{o_i\}_{i=1}^G \sim \pi_{\theta_{old}} (\cdot \mid q)$
      \State Evaluate rewards $\{r_i\}_{i=1}^{G}$ for all $o_i$ via $r_\varphi$
      \State Compute token-level advantages $\hat{A}_{i,t}$ for $t$-th token of $o_i$ using group-relative advantage estimation.
      \For{GRPO iteration = 1, \dots, $\mu$}
        \State Update policy model $\pi_\theta$ by maximizing the GRPO objective in Equation \ref{eq:GC-GRPO}
      \EndFor
    \EndFor 
    \State Refresh $r_\varphi$ using continual training with a replay buffer.
    \EndFor 
  \end{algorithmic}
  \textbf{Output} $\pi_\theta$
  \label{alg:iter-grpo}
\end{algorithm}

\subsubsection{Outcome Supervision RL with GRPO} More formally, for each prompt $q$, a collection of candidate outputs $\{o_1, o_2, \cdots, o_G\}$ is generated from the previous policy $\pi_{\theta_{old}}$. These candidates are assessed using a reward model, generating a set of $G$ rewards $\mathbf{r}=\{r_1, r_2, \cdots, r_G\}$. These rewards are standardized within the group by subtracting the mean and dividing by the standard deviation. In this outcome-based supervision, the terminal normalized reward for each output $o_i$ is treated as the advantage $\hat{A}_{i,t}$ for every token in $o_i$, i.e., $\hat{A}_{i, t} = \widetilde{r}_i = \frac{r_i- {\rm mean}(\mathbf{r})}{{\rm std}(\mathbf{r})}$. The policy is then updated to maximize the objective stated in equation (\ref{eq:GRPO-obj}).

\subsubsection{Process Supervision RL with GRPO} Relying solely on end-of-trajectory rewards—as in outcome supervision—may be inadequate or suboptimal for directing policy learning in intricate mathematical reasoning tasks. In line with \cite{wang2023math}, we also investigate process supervision, which attributes rewards to each intermediate reasoning step. Specifically, for question $q$ and group of outputs $\{o_1, o_2, \cdots, o_G\}$, a process reward model assigns step-wise rewards, yielding $\mathbf{R} = \{ \{r_1^{index(1)},\cdots,r_1^{index(K_1)}\}, \cdots,  \{r_G^{index(1)},\cdots,r_G^{index(K_G)}\} \}$, where $index(j)$ indicates the position of the terminal token of the $j$-th step and $K_i$ is the number of steps in output $o_i$. These step-level rewards are likewise standardized across the group: $\widetilde{r}_i^{index(j)} = \frac{r_i^{index(j)} - {\rm mean(\mathbf{R})}}{{\rm std(\mathbf{R})}}$. Subsequently, process supervision sets the advantage for each token as the sum of all subsequent normalized step rewards, i.e., $\hat{A}_{i, t} = \sum_{index(j) \ge t} \widetilde{r}_i^{index(j)}$, and the policy is then trained to maximize the objective described by equation (\ref{eq:GRPO-obj}).

\subsubsection{Iterative RL with GRPO} 

As reinforcement learning advances, the existing reward model may become inadequate for guiding the current policy model. To address this, we implement an iterative RL approach leveraging GRPO. As illustrated in Algorithm \ref{alg:iter-grpo}, this procedure involves periodically creating updated datasets for the reward model by sampling outputs from the latest policy model. The reward model is subsequently trained using a replay strategy that integrates 10\% of prior training samples. After this, the reference model is synchronized with the latest policy model, and the policy model itself undergoes further training, this time supervised by the newly improved reward model.

\subsection{Training and Evaluating DeepSeekMath-RL}

Our RL experiments utilize the DeepSeekMath-Instruct 7B model as a foundation. For RL training, we draw on chain-of-thought-style questions from GSM8K and MATH within the SFT dataset, totaling approximately 144K questions. Other SFT examples are omitted so as to assess the effects of RL on test sets that are not included during RL training. Construction of reward model training sets follows the methodology outlined in \citep{wang2023math}. The reward model is initially fine-tuned from DeepSeekMath-Base 7B, employing a learning rate of 2e-5. In the GRPO phase, the policy model’s learning rate is set to 1e-6, and a KL penalty coefficient of 0.04 is utilized. For every question, we generate $64$ response samples, with the maximum sequence length capped at 1024 and a batch size of 1024 used throughout training. The policy model receives just a single parameter update after each exploration cycle. For assessment, DeepSeekMath-RL 7B is tested on the same benchmarks as DeepSeekMath-Instruct 7B. Specifically, chain-of-thought reasoning benchmarks GSM8K and MATH are treated as in-domain, while all other tasks are designated as out-of-domain evaluations.

Table \ref{tab:sft_rl_math} presents the results achieved by both open-source and closed-source models employing chain-of-thought and tool-augmented reasoning on English and Chinese evaluation sets. Our observations include: 1) With chain-of-thought reasoning, DeepSeekMath-RL 7B obtains accuracies of 88.2\% on GSM8K and 51.7\% on MATH, outperforming every open-source model in the 7B–70B range as well as the bulk of closed-source models. 2) Notably, DeepSeekMath-RL 7B is trained solely using chain-of-thought-format instruction tuning data from GSM8K and MATH, with DeepSeekMath-Instruct 7B as its initialization. Although its training data is relatively limited in scope, DeepSeekMath-RL 7B consistently surpasses DeepSeekMath-Instruct 7B across evaluation criteria, providing evidence of the advantage brought by reinforcement learning.

\section{Discussion}
This section outlines insights obtained from our pre-training and RL studies.

\subsection{Lessons Learnt in Pre-Training}

We begin by summarizing our findings regarding pre-training. Unless explicitly stated, the experimental settings correspond to those described in Section \ref{sec:quality-policy}. Specifically, references to the DeepSeekMath Corpus pertain to an 89B-token version derived from the second round of dataset assembly.

\subsubsection{Code Training Benefits Mathematical Reasoning}

Although widely presumed, the idea that code training can enhance reasoning skills has not been fully substantiated. We seek to partially address this question within mathematical contexts: our results indicate that code training bolsters models’ capacity for mathematical reasoning, irrespective of the use of auxiliary tools.

To examine the effect of code pre-training on mathematical reasoning, we performed experiments involving both two-stage and one-stage training regimes:

\textbf{Two-Stage Training}

\begin{itemize}[topsep=0pt]
    \item \textbf{Code Training for 400B Tokens $\rightarrow$ Math Training for 150B Tokens}: DeepSeek-LLM 1.3B is first pre-trained with 400B tokens of code data and then further trained using 150B tokens drawn from mathematical content;
    \item \textbf{General Training for 400B Tokens $\rightarrow$ Math Training for 150B Tokens}: For comparison, a parallel setup is employed in which 400B tokens are sampled from a comprehensive general-domain corpus compiled by DeepSeek-AI, replacing the code data in the initial stage. This aims to assess the specific benefits that code tokens confer, relative to general-domain tokens, for downstream mathematical reasoning capabilities.
\end{itemize}

\textbf{One-Stage Training}

\begin{itemize}[topsep=0pt]
    \item \textbf{Math Training for 150B Tokens}: DeepSeek-LLM 1.3B undergoes training solely on 150B math-specific tokens;
    \item \textbf{Training on a mixture of 400B Code Tokens and 150B Math Tokens}: It is observed that performing math training after code training results in diminished performance on coding tasks. To address this, a unified training stage is considered, wherein both 400B code tokens and 150B math tokens are mixed together, with the goal of simultaneously enhancing mathematical reasoning and preserving coding performance, thereby combating catastrophic forgetting.
\end{itemize}

\paragraph{Results}
Tables \ref{tab:code-math} and \ref{tab:code-math-general-eval} present the model’s downstream results across these different training configurations.

Introducing code tokens in the training regimen substantially advances the model’s proficiency in program-assisted mathematical reasoning, regardless of whether training proceeds in two stages or as a unified stage. Evidence from Table \ref{tab:code-math} reveals that pre-training exclusively on code notably bolsters the ability to tackle GSM8K and MATH datasets via Python. A subsequent phase of math-focused training leads to additional performance gains. Moreover, integrating code and math tokens in a joint training stage effectively curbs catastrophic forgetting, an issue characteristic of staged training, and additionally reinforces both coding skills (Table \ref{tab:code-math-general-eval}) and programmatic mathematical reasoning (Table \ref{tab:code-math}).

Beyond its impact on tool-based reasoning, pre-training on code also aids mathematical problem solving in scenarios absent tool usage. Specifically, code pre-training in the initial stage yields notable progress, and further enhances the subsequent math training, ultimately delivering the highest measured accuracy. Conversely, merging code and math tokens in a single-stage curriculum appears to hinder mathematical reasoning when tools are not employed, potentially due to the limited capacity of DeepSeek-LLM 1.3B, which may be insufficient to optimally internalize the combined learning objectives of both domains.

\subsubsection{ArXiv Papers Appear Unhelpful for Enhancing Mathematical Reasoning}

While arXiv papers are often included as part of the math pre-training dataset \citep{minerva,gpt-f,llemma,mathpile}, systematic investigations into their effects on mathematical reasoning performance remain limited. Surprisingly, our empirical findings indicate that exposure to arXiv content yields little benefit for mathematical reasoning tasks. To examine this, we utilized DeepSeek-LLM 1.3B and DeepSeek-Coder-Base-v1.5 7B \citep{deepseek-coder}, and trained these models on different versions of arXiv data subjected to diverse preprocessing workflows:

\begin{itemize}[topsep=0pt]
    \item \textbf{MathPile} \citep{mathpile}: This collection contains 8.9B tokens, compiled with filtering heuristics and cleaning steps, where over 85\% of content originates from scientific arXiv submissions;
    \item \textbf{ArXiv-RedPajama} \citep{redpajama}: An extensive dataset of 28.0B tokens constructed from raw arXiv LaTeX files, in which non-essential components such as preambles, macros, comments, and bibliographies are excluded.
\end{itemize}

Our procedure involved individually training DeepSeek-LLM 1.3B on 150B tokens and DeepSeek-Coder-Base-v1.5 7B on 40B tokens for each corpus described above. The resulting models trained solely on arXiv data failed to achieve measurable gains—and in certain instances even exhibited a decline—in performance across a range of mathematical evaluation tasks, irrespective of problem difficulty. The tested benchmarks include datasets focusing on quantitative reasoning like GSM8K and MATH (see Table \ref{tab:arxiv-cot}), STEM-oriented multiple-choice tasks such as MMLU-STEM (see Table \ref{tab:arxiv-cot}), as well as datasets dedicated to formal mathematics, like miniF2F (see Table \ref{tab:arxiv_atp}).

It is important to note several caveats to these observations. Our investigation did not consider:

\begin{itemize}[topsep=0pt]
    \item The utility of arXiv data for narrower mathematical tasks omitted from our evaluation, such as converting formal mathematical content into informal explanations;
    \item Possible synergies when arXiv data are used in conjunction with additional training corpora;
    \item The potential for improvements that might be seen in much larger-scale models trained on arXiv data.
\end{itemize}

Consequently, further research is warranted to address these open questions.

\subsection{Perspectives on Reinforcement Learning}
\subsubsection{Moving Toward a Unified Framework} Here, we present a unified approach for analyzing and understanding a spectrum of training algorithms, encompassing SFT, RFT, DPO, PPO, and GRPO. We also empirically investigate the primary components of this paradigm. Generally, the gradient with respect to the parameters $\theta$ for any of these training techniques may be represented as:

\begin{equation}
    \nabla_{\theta}\mathcal{J}_{\textcolor{red}{\mathcal{A}}}(\theta) = \mathbb{E}[\underbrace{(q,o) \sim \textcolor{red}{\mathcal{D}}}_{Data \ Source}]\left( \frac{1}{|o|} \sum_{t=1}^{|o|}  \underbrace{GC_{{\mathcal{A}}}(q, o, t, \textcolor{red}{\pi_{{rf}}})}_{Gradient \ Coefficient}  \nabla_{\theta}\log \pi_{\theta}(o_t | q, o_{<t})\right).
\label{eq:objective}
\end{equation}

This formulation highlights three principal elements: 1) \textit{Data Source $\mathcal{D}$}, specifying the dataset from which training samples are drawn; 2) \textit{Reward Function $\pi_{{rf}}$}, providing the basis for the reward signal used in training; 3) \textit{Algorithm $\mathcal{A}$}, which maps the data and corresponding reward signals to a gradient coefficient $GC$ that modulates the strength of the gradient for reinforcement or penalization. The following canonical approaches are considered within this unified framework:

\begin{itemize}
    \item  \textbf{Supervised Fine-tuning (SFT)}: SFT adapts a pretrained model by training it on datasets curated by human annotators.
    \item \textbf{Rejection Sampling Fine-tuning (RFT)}: RFT enhances an SFT-trained model by continuing its training on responses sampled from itself, with these samples filtered for answer correctness with respect to the original SFT prompts.
    \item \textbf{Direct Preference Optimization (DPO)}: DPO incrementally tunes the SFT model using extended outputs generated by the SFT model and employs a pairwise DPO objective for learning from preferences.
    \item \textbf{Online Rejection Sampling Fine-tuning (Online RFT)}: Unlike standard RFT, Online RFT initializes the policy with the SFT model and further tunes it using new responses sampled on-the-fly from the evolving policy itself.
    \item \textbf{PPO/GRPO}: PPO/GRPO methods begin by seeding the policy model with the SFT parameters and subsequently reinforce it using samples generated from the current iteration of the policy model.
\end{itemize}

The elements constituting these approaches are outlined in Table \ref{tab:data_policy}. For an in-depth derivation, consult Appendix \ref{app:analysis-rl}.

\paragraph{Observation about Data Source} We classify the origin of data into two types: online sampling and offline sampling. Online sampling indicates that the dataset is produced from the ongoing policy model’s exploration during training, whereas offline sampling involves collecting data from outputs generated by the initial SFT model. The methods RFT and DPO utilize offline sampling, while Online RFT and GRPO adopt online sampling strategies.

Referring to Figure \ref{fig:rl-analysis}, it is evident that Online RFT yields superior performance over RFT across both evaluated benchmarks. Initially, Online RFT performs similarly to RFT, but as training advances, it establishes a distinct lead, underscoring the benefits of an online data collection approach. This phenomenon is logical: early in training, the actor’s outputs largely mirror those of the SFT model, with minimal deviations present in the collected data. Over time, however, the differences between data generated by the actor and the SFT model increase, and the advantages of leveraging real-time sampling become more pronounced.

\paragraph{Observation about Gradient Coefficient} The reward signal is mapped to a gradient coefficient, which in turn is used to update model parameters. In our experiments, we distinguish between two reward function types: ‘Rule’ and ‘Model.’ The Rule-based function evaluates response quality solely on answer correctness, while the Model-based function relies on a learned reward model, itself trained on judgments made by the rule-based function. Examining Equations \ref{eq:GC-RFT} and \ref{eq:GC-GRPO}, a primary distinction arises: GRPO modulates its gradient coefficient according to the specific reward value determined by the reward model, thereby enabling graded reinforcement or penalization depending on response quality. Conversely, Online RFT lacks this granularity, applying uniform positive reinforcement to all correct responses and not imposing any penalty on incorrect ones.

As depicted in Figure \ref{fig:rl-analysis}, GRPO demonstrates superior effectiveness over online RFT, underscoring the advantage of modifying positive and negative gradient coefficients. Moreover, GRPO+PS achieves better results than GRPO+OS, suggesting that utilizing more granular, step-specific gradient coefficients yields measurable benefits. Our investigation further extends to iterative RL; we perform two successive iterations in our experiments. The outcomes in Figure \ref{fig:iter-rl} reveal that this iterative RL approach brings substantial gains, particularly evident in the first cycle.

\subsubsection{Why RL Works?} This work employs reinforcement learning on a selected subset of instruction tuning data, leading to remarkable improvements over the baseline instruction-tuned model. To probe the underlying reasons for RL's effectiveness, we analyze Pass@K and Maj@K accuracies for both the Instruct and RL-optimized models across two evaluation benchmarks. As shown in Figure \ref{fig:maj-pass}, RL provides a notable increase in Maj@K, whereas Pass@K remains largely unaffected. This pattern suggests that RL does not necessarily enhance the model's inherent abilities but rather strengthens the reliability of correct outputs within the TopK selections. \textbf{In essence, the observed improvements are primarily due to raising the probability that the correct response appears in TopK, rather than an uplift in the model's foundational competence.} Likewise, \citep{wang2023making} identified a \textbf{misalignment problem} in reasoning tasks handled by SFT models, and demonstrated that targeted preference alignment strategies can remedy this issue and boost performance \citep{yuan2023rrhf,song2023preference,wang2023making}.

\subsubsection{How to Achieve More Effective RL?}
\label{sec:effec_rl}
Our findings affirm that RL exhibits strong performance in mathematical reasoning scenarios. Additionally, we introduce a unified framework that encompasses several pivotal training techniques, categorizing each as either a straightforward or approximate form of RL. According to Equation \ref{eq:objective}, this unified framework comprises three central elements: Data Source, Algorithm, and Reward Function. We outline possible directions for future research concerning each of these components.

\paragraph{Data Source} The data source forms the foundational input for all training frameworks. In RL settings, this specifically refers to unlabeled question sets and outputs produced via the policy model. This study restricts its data source to questions originating from the instruction tuning phase, combined with a basic nucleus sampling procedure for generating responses. This design choice might partly explain why our RL setup boosts Maj@K performance predominantly. Prospective work will extend RL applications to out-of-distribution queries and will incorporate \textbf{advanced sampling (decoding) strategies} inspired by tree-search algorithms \citep{yao2023tree}. Additionally, the integration of \textbf{efficient inference techniques} \citep{xia-etal-2023-speculative,leviathan2023fast,kwon2023efficient,xia2024unlocking}, which govern how efficiently the policy model explores the solution space, is likely to have a substantial impact on future improvements.

\paragraph{Algorithms} Algorithms utilize both data and reward signals to adjust the gradient coefficients used in updating model parameters. According to Equation \ref{eq:objective}, contemporary approaches overwhelmingly \textbf{TRUST} the reward function's signal to modulate the conditional probability for producing specific tokens. Nevertheless, this trust cannot guarantee that the reward signals are invariably accurate, particularly when confronting highly intricate tasks. For instance, even in the PRM800K datasets \citep{lightman2023let}, which are annotated by skilled professionals, around 20\% of the annotations have been found to be erroneous\footnote{\url{https://github.com/openai/prm800k/issues/12\#issuecomment-1728491852}}. Consequently, it is imperative to investigate reinforcement learning algorithms that demonstrate resilience in the face of noisy reward feedback. We anticipate that alignment strategies such as \textbf{WEAK-TO-STRONG} \citep{burns2023weak} will fundamentally transform how learning algorithms are designed and implemented.

\paragraph{Reward Function} The reward function constitutes the foundation of the training feedback in RL, where it typically takes the form of a neural reward model. We identify three primary directions for enhancing reward models: 1) \textbf{Improving the generalization capacity of reward models.} Effective reward models must generalize to unseen distributions and advanced generative outputs; without this ability, reinforcement learning risks merely stabilizing language model distributions without real advancement in core abilities; 2) \textbf{Incorporating the representation of uncertainty in the reward model.} The quantification of uncertainty could serve as a crucial intermediary between weaker reward models and algorithms built on weak-to-strong paradigms; 3) \textbf{Developing methods for rapid construction of high-quality process reward models} that can deliver nuanced feedback during reasoning tasks \citep{lightman2023let,wang2023math}.

\section{Conclusion, Limitation, and Future Work} 

In this work, we introduce DeepSeekMath, which sets a new standard among open-source models on the MATH benchmark for competition mathematics, rivaling the accuracy of proprietary systems. DeepSeekMath~builds upon DeepSeek-Coder-v1.5 7B, benefiting from extended continual training on 500B tokens, with 120B math-specific tokens largely obtained from Common Crawl serving as a major part of the dataset. Through a comprehensive ablation analysis, we observe that web pages are a valuable resource for obtaining high-quality mathematical content, while data from arXiv appears less impactful than initially anticipated. We propose Group Relative Policy Optimization (GRPO)—a modification of Proximal Policy Optimization (PPO)—which significantly enhances mathematical reasoning skills while also reducing memory requirements. Our experimental findings confirm that GRPO continues to deliver improvements even as DeepSeekMath-Instruct 7B attains high performance levels on evaluation tasks. Additionally, we outline a unified conceptual framework for various techniques and highlight promising avenues for developing more efficient reinforcement learning strategies.

Despite DeepSeekMath~demonstrating robust performance in quantitative reasoning, its abilities in geometry and formal theorem proving are not as advanced as those seen in closed-source counterparts. For example, during preliminary testing, the model struggled with triangle- and ellipse-related problems, which may point to biases in how data were selected during pre-training and fine-tuning stages. Moreover, due to its current scale, DeepSeekMath~lags behind GPT-4 in few-shot generalization: GPT-4 exhibits measurable gains with few-shot prompts, whereas DeepSeekMath~shows nearly identical outcomes on zero-shot and few-shot tasks. Moving forward, our plans include refining our curated data selection processes to yield a richer and higher-quality pre-training dataset. We will also pursue the research directions discussed in Section \ref{sec:effec_rl}, aiming to advance the efficiency of reinforcement learning methods for LLMs.

\end{document}